\chapter{整体猜想的证明}\label{ch:16}
% 原PDF第75页，第16章开始。
在\chref{15}的局部输入之后，本章按照现代抛物诱导输入\cite{CCF24}，证明整体范畴 Langlands 定理。我们省略通常的半密度扭转，仍用 $\DMod(\Bun_G)$ 表示上文的自守范畴。
\begin{theorem*}{}
Langlands 函子给出等价
\[\mathbb L_G:\DMod(\Bun_G)\xrightarrow{\sim}\IndCoh_{\mathrm{Nilp}}(\LocSys_{\checkG}),\]
并且与导出 Satake、Eisenstein 级数、常数项函子和对偶相容。
\end{theorem*}
证明先处理抛物部分，即 Eisenstein 部分，再证明剩余的尖点等价。

\section{粗 Langlands 函子}\label{sec:16.1}\index{Langlands 函子!粗形式|hyperpage}\index{Langlands 函子!重整化|hyperpage}
\srcpage{76}
第一个构造通过\cite{GR24a}中的 Langlands 函子构造，产生其粗形式
\[\mathbb L_G^0:\DMod(\Bun_G)\longrightarrow\QCoh(\LocSys_{\checkG}).\]
它由广义消失定理刻画：$\Rep(\checkG)_{\mathrm{Ran}}$ 在 $\DMod(\Bun_G)$ 上的局部 Hecke 作用，经由 $\QCoh(\LocSys_{\checkG})$ 的谱作用分解。采用第\ref{ch:13}–\ref{ch:15}章的记号，Whittaker 系数函子与谱作用给出图表
\origdiagram{p076-coarse.pdf}{0.4485}
这里 $\operatorname{CS}$ 表示几何 Casselman–Shalika 等价。完整的 Langlands 函子是 $\mathbb L_G^0$ 通过遗忘函子

\begin{transnote}{含糊点拨}
$\mathbb L_G$ 是粗函子 $\mathbb L_G^0$ 穿过遗忘函子 $\IndCoh_{\mathrm{Nilp}}\to\QCoh$ 的提升：目标必须换成重整化版本 $\QCoh^{!}$（ind-完备化），因为在紧对象上只能先落进下有界凝聚子范畴，普通 $\QCoh$ 不够，于是函数只在紧对象上与 $\mathbb L_G^0$ 相符。[依据：原书 p.76 本段；并见 [GR24a]。]
\end{transnote}
\[\Psi:\IndCoh_{\mathrm{Nilp}}(\LocSys_{\checkG})\longrightarrow\QCoh(\LocSys_{\checkG})\]
的提升，即 $\mathbb L_G^0=\Psi\circ\mathbb L_G$。在紧对象上，首先落入一个下有界的凝聚子范畴；可呈示函子 $\mathbb L_G$ 通过 ind-完备化得到。方便起见，用 $\QCoh^!(\LocSys_{\checkG})$ 表示相应的重整化谱目标。因此，简记为
\[\DMod(\Bun_G)\xrightarrow{\mathbb L_G}\QCoh^!(\LocSys_{\checkG}).\]
与导出 Satake 的相容性内置于这一构造中。新的工作是与 Eisenstein 和常数项函子的相容性\cite{CCF24}。

\section{Eisenstein 与常数项的相容性}\label{sec:16.2}\index{Eisenstein 函子!Langlands 相容性|hyperpage}\index{常数项函子!Langlands 相容性|hyperpage}
令 $P\subset G$ 为 Levi 商为 $M$ 的抛物子群，并令 $\check P\subset\checkG$ 为对偶抛物子群。忽略标准的移位与扭转，自守 Eisenstein 函子和常数项函子为
\[\operatorname{Eis}_P:\DMod(\Bun_M)\longrightarrow\DMod(\Bun_G),\qquad
\CT_P:\DMod(\Bun_G)\longrightarrow\DMod(\Bun_M).\]
在谱侧，它们由以下对应定义：
\[\LocSys_{\check M}\xleftarrow{q}\LocSys_{\check P}\xrightarrow{p}\LocSys_{\checkG}.\]
目标陈述是以下两个图表交换：
\origdiagramtag{p076-eis.pdf}{0.4645}{16.1}{eq:16.1}
以及
\origdiagramtag{p076-ct.pdf}{0.4713}{16.2}{eq:16.2}
\srcpage{77}
按\cite{CCF24}的形式，\eqref{eq:16.1}的证明归约到两个方块：一个涉及 Whittaker 系数的几何方块，以及一个涉及表示的谱方块。在几何 Casselman–Shalika 等价下，$\Whit(H)_x\simeq\Rep(\check H)_x$，所以这一比较变为 Whittaker 系数的抛物限制与表示上的相应 Chevalley 函子之间的相容性。

\section{谱方块：推前与局部化}\label{sec:16.3}\index{Chevalley 函子|hyperpage}\index{Ran 空间!标记点插入|hyperpage}
基本问题如下。对同态 $H_1\to H_2$，有映射 $\LocSys_{H_1}\to\LocSys_{H_2}$；对点 $x\in X$，还有局部的局部化函子
\[\operatorname{Loc}_x:\Rep(H_i)_x\longrightarrow\QCoh^!(\LocSys_{H_i}).\]
这里不能朴素地用一个方块，将沿 $\LocSys_{H_1}\to\LocSys_{H_2}$ 的拉回与表示限制比较。正确的方块使用推前和 Ran 插入。
\begin{lemma}{}{16.3.1}
假设 $H_1\to H_2$ 的核为幺幂群。令 $\operatorname{Ran}_{x,c}:=\{I\in\operatorname{Ran}(X)\mid x\in I\}$。则下面的图表交换：
\origdiagramtag{p077-general.pdf}{0.4713}{16.3}{eq:16.3}
\end{lemma}

这里，$\operatorname{inv}^+$ 先将对象插入标记点，再对幺幂核取导出不变量，等价地，取 Chevalley 上链。具体地，若 $V_x\in\Rep(H_1)_x$，则在点 $(x,y_1,\ldots,y_m)\in\operatorname{Ran}_{x,c}$，插入为

\begin{transnote}{含糊点拨}
先插入再取不变量，是为了与 Ran 空间的因子化结构相容：插入点 $x$ 上多出一个因子，再取幺幂核的导出不变量（即 Chevalley 上链 $C^{\bullet}(\check{\mathfrak u},-)$），剩下的正好是 Whittaker 归一化。每个插入点各带一个这样的因子，所以后文的公式是逐点乘起来的，而不是全局取一次。[依据：原书 p.77 本段；并见 [CCF24]。]
\end{transnote}
\[\operatorname{ins}(V_x)_{(x,y_1,\ldots,y_m)}=V_x\otimes\mathbf1_{y_1}\otimes\cdots\otimes\mathbf1_{y_m}.\]
\begin{proof}
证明是以下初等原理的 Ran 空间形式。若 $\pi:Y\to X$ 是一个 $D$-概形，且 $Y$ 拟仿射、余仿射，或满足相同的下降条件，则
\[\QCoh(Y)\simeq\pi_*\cO_Y\text{-}\mathrm{Mod}(\QCoh(X)),\]
其中 $\pi_*\cO_Y$ 一般是一个复形。例如，$\mathbb A^2\setminus\{0\}\to\mathrm{pt}$ 和 $B\Ga\to\mathrm{pt}$ 说明，必须保留这个复形，而不能只保留它的零次函数。将这一原理应用于形式圆盘，并使用插入与 Ran 卷积的相容性，就得到\eqref{eq:16.3}。
\end{proof}

对抛物商 $\check P\to\check M$，它的核是幺幂根基 $\check U$，引理给出
\origdiagram{p077-parabolic.pdf}{0.4690}
\srcpage{78}
再与 $\LocSys_{\check P}\to\LocSys_{\checkG}$ 复合，便得到常数项方向的谱方块
\origdiagram{p078-ct.pdf}{0.4690}
这里 $\operatorname{Chev}^+$ 是 Chevalley 函子。若 $V_x\in\Rep(\checkG)_x$，则示意地，
\[\operatorname{Chev}^+(V_x)_{(x,y_1,\ldots,y_m)}=
C^\bullet(\check{\mathfrak u},V_x)_x\otimes C^\bullet(\check{\mathfrak u})_{y_1}\otimes\cdots\otimes C^\bullet(\check{\mathfrak u})_{y_m},\]
其中 $\check{\mathfrak u}=\operatorname{Lie}(\check U)$ 是 $\check P$ 的幺幂根基的 Lie 代数。取对偶之后，等价地，在表示侧转到左伴随之后，出现相反的幺幂根基，给出 Eisenstein 方向的方块
\origdiagram{p078-eis.pdf}{0.4781}

\section{几何方块与半无限范畴}\label{sec:16.4}\index{半无限范畴|hyperpage}\index{第二伴随性|hyperpage}
几何侧是下面的方块：
\origdiagram{p078-geometric.pdf}{0.4485}
这里，$J^-:\Whit(M)_{\mathrm{Ran}_{x,c}}\to\Whit(G)_x$ 是 Eisenstein 方块中诱导方向的函子。同一个半无限核先给出局部限制函子 $r_P^-:\Whit(G)\to\Whit(M)$；正如第二伴随性所预言的，$J^-$ 是它的伴随形式。
\Needspace{4\baselineskip}
在局部，识别
\[\Whit(G)=\DMod(\Gr_G)^{LN,\chi_G},\qquad
\Whit(M)=\DMod(\Gr_M)^{LN_M,\chi_M}.\]
核从负半无限范畴构造：
\[\begin{aligned}
\operatorname{SI}_P^-&:=\DMod(\Gr_G)^{LN^-\cdot L^+M}\\
&\simeq(\DMod(\Gr_G)\otimes\DMod(\Gr_M))^{LP^-}\\
&\simeq\operatorname{Fun}^{LP^-}(\DMod(\Gr_G),\DMod(\Gr_M)).
\end{aligned}\]
因此，任意 $LP^-$-线性函子 $F:\DMod(\Gr_G)\to\DMod(\Gr_M)$ 都给出一个限制型函子
\[r_F:\DMod(\Gr_G)^{LN,\chi_G}\longrightarrow\DMod(\Gr_M)^{LN_M,\chi_M}.\]
上面使用的函子 $J^-$，是与 $\operatorname{SI}_P^-$ 的单位对象对应的伴随，即与轨道
\[LN^-\cdot1\cdot L^+G/L^+G\subset\Gr_G\]
上的标准对象对应的伴随。
\srcpage{79}
然后，将 Whittaker 特征标 $\chi_G$ 退化为一个在 $P$ 的幺幂根基上平凡、并在 $LN_M$ 上诱导 $\chi_M$ 的特征标 $\chi_P$。这给出局部模式
\[\DMod(\Gr_G)^{LN,\chi_G}\xrightarrow{\text{邻近圈}}\DMod(\Gr_G)^{LN,\chi_P}\longrightarrow\DMod(\Gr_M)^{LN_M,\chi_M}.\]
半无限对应使用 $P^-$，而特征标退化使用 $P$，这是第二伴随性的一个表现。

令 $\operatorname{Chev}^-:=(\operatorname{Chev}^+)^L$。在 Casselman–Shalika 等价下，所需的顶面方块如下：
\origdiagram{p079-chevalley.pdf}{0.4440}
证明使用 $\operatorname{SI}_P^-$ 的谱描述：
\begin{equation}
\begin{split}
\operatorname{SI}_P^-\simeq\QCoh^!\bigl(
\LocSys_{\checkG}(D)\times_{\LocSys_{\checkG}(D^\circ)}\LocSys_{\check P}(D^\circ)\\
\times_{\LocSys_{\check M}(D^\circ)}\LocSys_{\check M}(D)\bigr).
\end{split}
\tag{16.4}\label{eq:16.4}
\end{equation}

\begin{transnote}{含糊点拨}
四个因子依次读：$D$ 上的整体局部系统、它到 $D^\circ$ 的限制、$P$ 的抛物约化、$M$ 侧的局部系统。$P=G$ 时 $\LocSys_{\check P}=\LocSys_{\check G}$，右端退化为导出 Satake 的谱侧——这正是检验半无限范畴与谱侧对应是否正确的边界情形。[依据：原书 p.79 本式；并见 [CCF24]。]
\end{transnote}
当 $P=G$ 时，这恰好是导出 Satake 的谱侧。更一般地，它是导出 Satake 的抛物半无限对应版本。

局部 Langlands 的直观思想是：一个 $LG$-模范畴应当对应于 $\LocSys_{\checkG}(D^\circ)$ 上的范畴层。将环路群从 $G$ 限制到 $M$，对应于沿 $\LocSys_{\check M}(D^\circ)\to\LocSys_{\checkG}(D^\circ)$ 拉回。因此，预期有
\[\begin{aligned}
\DMod(\Gr_G)&\rightsquigarrow\QCoh(\LocSys_{\checkG}(D)),\\
\DMod(\Gr_M)&\rightsquigarrow\QCoh(\LocSys_{\check M}(D)),\\
\DMod(\Gr_G)^{LM}&\rightsquigarrow\QCoh\bigl(\LocSys_{\checkG}(D)\times_{\LocSys_{\checkG}(D^\circ)}\LocSys_{\check M}(D^\circ)\bigr).
\end{aligned}\]
在两侧计算 $\operatorname{Fun}(\text{--},\text{--})$，就得到\eqref{eq:16.4}，从而证明 Eisenstein 相容性\eqref{eq:16.1}。

\section{为何常数项更困难}\label{sec:16.5}\index{Kac–Moody–Whittaker 对偶|hyperpage}
证明\eqref{eq:16.1}的方法不能直接证明常数项相容性\eqref{eq:16.2}。在 Eisenstein 图表中，所有箭头都指向推前所需的方向。对于 $\CT_P$，必须反转这些箭头，而简单的几何论证不再适用。

替代方法是 Kac–Moody–Whittaker 对偶。在非临界水平，比较
\origdiagram{p079-noncritical.pdf}{0.4599}
特殊化到临界水平，等价地，特殊化到 FLE 退化中量子参数的边界值后，局部 Whittaker 范畴被球面临界 Kac–Moody 范畴替换。在 $\IndCoh$ 约定下，FLE 将它们识别为非分歧 oper：
\[\operatorname{KL}_{G,\mathrm{crit}}\simeq\IndCoh(\Op^{\mathrm{unr}}_{\checkG}),\qquad
\operatorname{KL}_{M,\mathrm{crit}}\simeq\IndCoh(\Op^{\mathrm{unr}}_{\check M}).\]
相应的局部函子为
\[r_P^{\mathrm{KM}}:\operatorname{KL}_{G,\mathrm{crit}}\longrightarrow\operatorname{KL}_{M,\mathrm{crit}},\qquad
M\longmapsto C^{\infty/2}(\mathfrak u((t)),M),\]
\srcpage{80}
其中 $\mathfrak u$ 是 $P$ 的幺幂根基的 Lie 代数。它的谱对应物是函子
\[r_{\check P}^{\mathrm{op}}:\IndCoh(\Op^{\mathrm{unr}}_{\checkG})\longrightarrow\IndCoh(\Op^{\mathrm{unr}}_{\check M}),\]
由 $\operatorname{SI}_{\check P^-}^{\mathrm{spec}}$ 中的单位对象给出。这些函子证明了常数项立方体的顶面与侧面\cite{CCF24}。剩余的前面与后面并非仅由半无限论证推出，而是通过手征方法证明。这就完成了常数项相容性\eqref{eq:16.2}的证明。

\section{伴随与 Eisenstein 生成的部分}\label{sec:16.6}\index{Eisenstein 生成子范畴|hyperpage}\index{可约局部系统轨迹|hyperpage}\index{不可约局部系统轨迹|hyperpage}
Eisenstein 论证的最后部分解释，相容性\eqref{eq:16.1}和\eqref{eq:16.2}如何进入整体归纳。令 $\DMod(\Bun_G)_{\mathrm{Eis}}$ 为由真抛物子群 $P$ 对应的 $\operatorname{Eis}_P(\DMod(\Bun_M))$ 生成的子范畴，令 $\LocSys^{\mathrm{red}}_{\checkG}\subset\LocSys_{\checkG}$ 为可约轨迹，即所有真抛物子群 $\check P$ 对应的 $\LocSys_{\check P}\to\LocSys_{\checkG}$ 的像的并。等价地，这些局部系统容许真抛物约化。其补集是不可约轨迹 $\LocSys^{\mathrm{irr}}_{\checkG}$。

需要证明粗函子 $\mathbb L_G^0$ 有左伴随。要点是，只需在生成对象上构造左伴随。在谱侧，生成对象来自 $\operatorname{Eis}^{\mathrm{spec}}_{\check P}(\QCoh(\LocSys_{\check M}))$，而 Eisenstein 相容性通过从 $M$ 输运伴随，给出所需的自守函子。等价地，使用以下事实：转到 Ran 作用描述后，粗函子已经有左伴随；$\Rep(\check M)_{\mathrm{Ran}}$ 作用于自守侧，而商 $\Rep(\check M)_{\mathrm{Ran}}\to\QCoh(\LocSys_{\check M})$ 识别谱作用。

因此，假设每个真 Levi 子群 $M$ 的 Langlands 等价已经成立，则函子 $\mathbb L_G$ 在 Eisenstein 直和因子上诱导等价\cite{CCF24}：
\begin{equation}
\mathbb L_G:\DMod(\Bun_G)_{\mathrm{Eis}}\xrightarrow{\sim}\QCoh^!(\LocSys_{\checkG})_{\mathrm{red}}.
\tag{16.5}\label{eq:16.5}
\end{equation}
于是，整体猜想被归约到尖点部分。剩余的陈述是
\[\mathbb L_{G,\mathrm{cusp}}:\DMod(\Bun_G)_{\mathrm{cusp}}\xrightarrow{\sim}\QCoh(\LocSys^{\mathrm{irr}}_{\checkG}).\]
这个归约为下面的尖点论证提供输入。

现在，在上述 Eisenstein 归约之后，按照证明系列中的重数一定理部分\cite{GR24b}，完成证明。仍用记号
\[\mathbb L_G:\DMod(\Bun_G)\longrightarrow\QCoh^!(\LocSys_{\checkG})\]
表示重整化 Langlands 函子，并用 $\mathbb L_G^L$ 和 $\mathbb L_G^R$ 表示它在相关直和因子上限制后的左伴随和右伴随。前面各节构造了 $\mathbb L_G$、构造了 $\mathbb L_G^L$，并证明 Eisenstein 等价
\[\mathbb L_G:\DMod(\Bun_G)_{\mathrm{Eis}}\xrightarrow{\sim}\QCoh^!(\LocSys_{\checkG})_{\mathrm{red}}.\]
剩下的任务是证明尖点等价
\[\mathbb L_{G,\mathrm{cusp}}:\DMod(\Bun_G)_{\mathrm{cusp}}\xrightarrow{\sim}\QCoh(\LocSys^{\mathrm{irr}}_{\checkG}).\]
这里，$\LocSys^{\mathrm{irr}}_{\checkG}$ 是不可约局部系统的开子叠，即不容许约化到真抛物子群 $\check P\subsetneq\checkG$ 的局部系统。

\section{尖点函子的已知输入}\label{sec:16.7}\index{保守函子!尖点 Langlands 函子|hyperpage}\index{单子|hyperpage}
\srcpage{81}
第一个输入是保守性。
\begin{theorem}{}{16.7.1}
函子 $\mathbb L_{G,\mathrm{cusp}}$ 是保守的。
\end{theorem}
\begin{proof}
证明思路大致使用 Whittaker 系数与几何 Ramanujan 输入\cite{FR25,Ber21,GR24b}。自守系数函子给出如下形式的方块：
\origdiagram{p081-cuspidal.pdf}{0.4918}
下方的水平函子是 Casselman–Shalika 识别。由几何 Ramanujan 输入，尖点对象是温和的，而 Whittaker 系数函子在温和子范畴上保守。因此，$\mathbb L_{G,\mathrm{cusp}}$ 是保守的。
\end{proof}

第二个输入是\secref{16.6}中已经构造的伴随的存在性。因此，要证明 $\mathbb L_{G,\mathrm{cusp}}$ 是等价，只需识别单子
\[\mathbb L_{G,\mathrm{cusp}}\circ\mathbb L_{G,\mathrm{cusp}}^L\in\End(\QCoh(\LocSys^{\mathrm{irr}}_{\checkG})).\]
谱范畴是 $\QCoh(\LocSys^{\mathrm{irr}}_{\checkG})$-线性的。因此，这个单子是与结合代数对象
\[\begin{split}
\mathcal A_{G,\mathrm{irr}}&:=(\mathbb L_{G,\mathrm{cusp}}\circ\mathbb L_{G,\mathrm{cusp}}^L)(\cO_{\LocSys^{\mathrm{irr}}_{\checkG}})\\
&\in\operatorname{AssocAlg}(\QCoh(\LocSys^{\mathrm{irr}}_{\checkG}))
\end{split}\]
作张量积。等价地，若在整个谱侧取 $\mathcal A_G=\mathbb L_G\mathbb L_G^L(\cO)$，则 $\mathcal A_{G,\mathrm{irr}}\simeq\mathcal A_G|_{\LocSys^{\mathrm{irr}}_{\checkG}}$。这一限制并非仅由 $\QCoh^!$ 的记号形式地推出，而是来自重整化谱范畴与不可约开集上的普通 $\QCoh$-模结构的相容性。因此，唯一剩下的任务是
\[\cO_{\LocSys^{\mathrm{irr}}_{\checkG}}\xrightarrow{\sim}\mathcal A_{G,\mathrm{irr}}.\]
这个陈述的可约部分，已在上文通过 Eisenstein 归纳证明。

\section{双伴随性与自对偶代数}\label{sec:16.8}\index{双伴随性|hyperpage}\index{自对偶代数|hyperpage}
\begin{theorem}{双伴随性\cite{ABC24b}}{16.8.1}
存在伴随的典范等价
\[\mathbb L_{G,\mathrm{cusp}}^L\simeq\mathbb L_{G,\mathrm{cusp}}^R.\]
\end{theorem}
\begin{proof}
证明概述如下。通过上面使用的局部化与系数图表比较两侧。自守尖点范畴和 $\QCoh(\LocSys^{\mathrm{irr}}_{\checkG})$ 都带有自然对偶。$\mathbb L_G$ 与 FLE、谱 Poincaré 函子和 Whittaker 系数的相容性，将 $\mathbb L_{G,\mathrm{cusp}}$ 的对偶函子识别为同一个函子。因此，左伴随和右伴随相同。
\end{proof}

其结果是：代数 $\mathcal A_{G,\mathrm{irr}}$ 作为 $\LocSys^{\mathrm{irr}}_{\checkG}$ 上的复形是完美的，并且典范自对偶。事实上，右伴随给出与左伴随相同的核，但采用相反的伴随约定。因此，与 $\mathcal A_{G,\mathrm{irr}}$ 对偶的对象是同一个对象，只是现在将它看作余代数。将这一余代数记为 $\mathcal B_{G,\mathrm{irr}}$。用公式表示，
\[\mathcal B_{G,\mathrm{irr}}:=\mathcal A_{G,\mathrm{irr}}^\vee,\qquad
\mathcal A_{G,\mathrm{irr}}^\vee\simeq\mathcal A_{G,\mathrm{irr}}.\]
\srcpage{82}
正是在这一步，双伴随性将单子问题转化为关于谱侧伪真 Ran 推前的几何陈述。

\section{一般位置 oper 与余代数}\label{sec:16.9}\index{oper!扩展一般位置结构|hyperpage}\index{余代数|hyperpage}\index{有限 étale 代数|hyperpage}
令
\[\pi:\Op^{\mathrm{gen}}_{\checkG,\mathrm{irr}}\longrightarrow\LocSys^{\mathrm{irr}}_{\checkG}\]
为不可约局部系统上的一般位置，更准确地说是扩展一般位置 $\checkG$-oper 结构的叠。非正式地说，在 $\sigma\in\LocSys^{\mathrm{irr}}_{\checkG}$ 上的纤维，是 $\sigma$ 到 $\check B$ 的一般位置约化所组成的空间，这些约化在 $X$ 的有限多个点之外满足 oper 横截条件。FLE／局部化形式体系将上述余代数识别为 de Rham 直接像\cite{ABC24a,GR24b}
\[\mathcal B_{G,\mathrm{irr}}\simeq\pi_{\mathrm{dR},!}\omega_{\Op^{\mathrm{gen}}_{\checkG,\mathrm{irr}}},\]
这是 $\QCoh(\LocSys^{\mathrm{irr}}_{\checkG})$ 中的余交换余代数对象。在证明中，此对象由真空 Poincaré 核构造：局部 FLE 说明，Kac–Moody 真空对象被送到非分歧 oper 对象，而 Ran 因子化将局部的一般位置约化粘合为上述全局一般位置 oper 对应。

下面使用几个初等推论。第一，$\mathcal B_{G,\mathrm{irr}}$ 余交换，位于上同调次数 $\le0$，并沿 $\LocSys^{\mathrm{irr}}_{\checkG}$ 带有平坦联络。第二，由自对偶性，$\mathcal A_{G,\mathrm{irr}}$ 交换，位于次数 $\ge0$。由于对偶将这一完美对象送回自身，反向的次数界也成立；所以 $\mathcal A_{G,\mathrm{irr}}$ 位于心范畴中，并且是带有平坦联络的向量丛。第三，$\pi_{\mathrm{dR},!}\omega$ 具有有限单值性。因此：
\begin{theorem}{}{16.9.1}
代数 $\mathcal A_{G,\mathrm{irr}}$ 是有限 étale 交换 $\cO_{\LocSys^{\mathrm{irr}}_{\checkG}}$-代数。等价地，它定义 $\LocSys^{\mathrm{irr}}_{\checkG}$ 的有限 étale 覆盖。
\end{theorem}
\begin{note}\numberednote{16.9.2}
对经典群，Beraldo–Kazhdan–Schlank 的定理加强了这个有限单值性输入：相关的一般位置 oper 纤维是同调可缩的\cite{BKS18}。因此，在 $A$、$B$、$C$、$D$ 型中，该有限 étale 覆盖已被迫为平凡覆盖。证明以此为模型情形，再将一般情形归约到同样的有限 étale 论证。
\end{note}

\section{归约到一个数值秩}\label{sec:16.10}\index{重数一定理|hyperpage}\index{Cohen–Macaulay 性质|hyperpage}
现在限制到证明中使用的标准良好情形：$G$ 为伴随型，且 $g_X>1$；小的例外情形 $g_X=2$ 且 $\checkG=\SL_2$ 通过单独的归约处理。在这一情形中，$\LocSys^{\mathrm{irr}}_{\checkG}$ 在这里所需的意义下没有非平凡的有限 étale 覆盖。因此，它上面的任意有限 étale 代数都是常值的：
\[\mathcal A_{G,\mathrm{irr}}\simeq\cO_{\LocSys^{\mathrm{irr}}_{\checkG}}^{\oplus n},\]
其中 $n\ge1$ 为整数。要证明 $n=1$，只需证明全局截面上的映射
\[\Gamma(\LocSys^{\mathrm{irr}}_{\checkG},\cO)\longrightarrow\Gamma(\LocSys^{\mathrm{irr}}_{\checkG},\mathcal A_{G,\mathrm{irr}})\]
是同构。在良好情形中，不可约轨迹的补集余维至少为二，而且局部系统叠是 Cohen–Macaulay 的；因此，将 $\mathcal A_{G,\mathrm{irr}}$ 延拓为整个 $\LocSys_{\checkG}$ 上的全局代数 $\mathcal A_G$ 后，可以检验同一个问题。只需证明
\[\Gamma(\LocSys_{\checkG},\cO)\xrightarrow{\sim}\Gamma(\LocSys_{\checkG},\mathcal A_G).\]
\srcpage{83}
全局函数方面的输入给出 $\Gamma(\LocSys_{\checkG},\cO)=k$。由全局 Poincaré 真空核的定义，右端是基本 Whittaker 对象 $W_i$ 的自同态代数：
\[\Gamma(\LocSys_{\checkG},\mathcal A_G)\simeq\End(W_i).\]
因此，最后需要计算的是 $\End(W_i)=k$。

计算中所用的 Whittaker 图卡，是通常的 Drinfeld 紧化中的开分层。对充分负的次数数据 $i$，有一个局部闭图：
\origdiagram{p083-whittaker.pdf}{0.4850}
其中 $r(G)=\operatorname{rank}_{\mathrm{ss}}(G)$，$\chi_{\mathrm{glob}}$ 是简单根 Whittaker 特征标的乘积。相关的极小层形如
\[\mathring{\mathcal A}_{\min}\simeq\Gm^{r(G)}\times BU_i,\]
其中 $U_i$ 为幺幂群。对于洁净 Whittaker 对象，$BU_i$ 因子只贡献单位对象；因此，相关的自同态计算归约为
\[\End_{\DMod(\mathring{\mathcal A}_{\min})}\bigl(W_i|_{\mathring{\mathcal A}_{\min}}\bigr)
\simeq\End_{\DMod(\Gm^{r(G)})}\bigl(\mathcal E^{\boxtimes r(G)}\bigr).\]
对象 $W_i$ 在这一图卡上限制为洁净指数对象 $\chi_{\mathrm{glob}}^!\mathcal E$。因此，这一计算就是 $\Gm^{r(G)}$ 上 $\mathcal E^{\boxtimes r(G)}$ 的初等计算，并且
\[\End(W_i)=k.\]
所以 $n=1$，从而 $\mathcal A_{G,\mathrm{irr}}\simeq\cO$。于是，单子 $\mathbb L_{G,\mathrm{cusp}}\mathbb L_{G,\mathrm{cusp}}^L$ 是恒等单子。由于 $\mathbb L_{G,\mathrm{cusp}}$ 是保守的，并且具有所需的伴随，它是一个等价：
\[\mathbb L_{G,\mathrm{cusp}}:\DMod(\Bun_G)_{\mathrm{cusp}}\xrightarrow{\sim}\QCoh(\LocSys^{\mathrm{irr}}_{\checkG}).\]
再与可约部分的 Eisenstein 等价\eqref{eq:16.5}结合，便证明了整体几何 Langlands 等价\cite{CCF24,ABC24b,GR24b}。
